\section{Discrete Numerical Architecture and Stability Proofs}
\label{sec:discrete_proofs}

To translate the continuous geometric and dynamical formulations of GSPT into discrete atmospheric models without introducing unphysical grid-point coalescence or explicit time-step crashes, we establish an adaptive grid-stretching algorithm alongside an invariant manifold solver \cite{nieuwstadt1984}.

\subsection{Discrete Equidistribution and Bounded Grid Stretching}

Let $z \in [0, H_{\text{SBL}}]$ represent physical vertical height, and let $\zeta_{zz}(z)$ denote the vertical coordinate curvature under local Nieuwstadt scaling. We define the adaptive physical coordinate density function $\rho_{\text{clamp}}(z)$ as:
\begin{equation}
    \rho_{\text{clamp}}(z) \equiv \max\left( \frac{k}{\Delta z_{\text{max}}}, \; \min\left( \frac{k}{\Delta z_{\text{min}}}, \; \left| \zeta_{zz}(z) \right| + \epsilon_z \right) \right), \label{eq:clamped_density}
\end{equation}
where $\Delta z_{\text{min}}$ and $\Delta z_{\text{max}}$ represent the minimum and maximum vertical grid steps, $\epsilon_z > 0$ is a coordinate regularization parameter, and $k > 0$ is a normalization scalar. Let $C(z) = \int_0^z \rho_{\text{clamp}}(s) \, ds$ define the cumulative density function. The physical vertical levels $z_j$ (for $j = 1, \dots, N_z$) are generated by uniform equidistribution of the computational coordinate $y \in [0, 1]$:
\begin{equation}
    y_j = \frac{j - 1}{N_z - 1}, \qquad z_j = y^{-1}(y_j). \label{eq:node_mapping}
\end{equation}

\begin{theorem}[Discrete Grid-Spacing Bounds]
\label{thm:grid_bounds}
Under the bounded equidistribution density \eqref{eq:clamped_density}, the discrete physical grid steps $\Delta z_j \equiv z_j - z_{j-1}$ are globally bounded by:
\begin{equation}
    \Delta z_{\text{min}} \le \Delta z_j \le \Delta z_{\text{max}} \qquad \forall j \in \{2, \dots, N_z\}. \label{eq:discrete_bounds_theorem}
\end{equation}
\end{theorem}

\begin{proof}
Applying the Mean Value Theorem to the normalized cumulative density function $y(z)$ on the interval $[z_{j-1}, z_j]$ guarantees an intermediate point $\xi \in (z_{j-1}, z_j)$ such that:
\begin{equation}
    y(z_j) - y(z_{j-1}) = y'(\xi) (z_j - z_{j-1}) \quad \Longrightarrow \quad \frac{1}{N_z - 1} = \frac{\rho_{\text{clamp}}(\xi)}{C_{\text{max}}} \Delta z_j. \label{eq:mvt_proof_step}
\end{equation}
Solving for the local grid step yields $\Delta z_j = \frac{C_{\text{max}}}{(N_z - 1) \rho_{\text{clamp}}(\xi)}$. By definition, the clamped density is bounded pointwise by $\frac{k}{\Delta z_{\text{max}}} \le \rho_{\text{clamp}}(z) \le \frac{k}{\Delta z_{\text{min}}}$. Integrating this across the compact domain yields the bounds on the cumulative mass:
\begin{equation}
    \frac{k H_{\text{SBL}}}{\Delta z_{\text{max}}} \le C_{\text{max}} \le \frac{k H_{\text{SBL}}}{\Delta z_{\text{min}}}. \label{eq:cmax_bounds}
\end{equation}
Setting the computational scaling factor to $k \equiv \frac{C_{\text{max}}}{N_z - 1}$, the grid-spacing relation simplifies to $\Delta z_j = \frac{k}{\rho_{\text{clamp}}(\xi)}$. Substituting the pointwise density bounds into this expression yields:
\begin{equation}
    \Delta z_{\text{min}} \le \Delta z_j \le \Delta z_{\text{max}}, \label{eq:bounds_resolved}
\end{equation}
completing the proof.
\end{proof}

\begin{theorem}[Asymptotic Grid Convergence]
\label{thm:grid_convergence}
Let $\chi(z) \in C^2$ represent a smooth, twice-differentiable inverse Obukhov length profile. Then, the discrete vertical coordinates $z_j$ converge asymptotically to the continuous optimal coordinate path $z(y)$ with quadratic accuracy in the vertical level resolution:
\begin{equation}
    \max_{1 \le j \le N_z} \left| z_j - z(y_j) \right| = \mathcal{O}\left( N_z^{-2} \right). \label{eq:convergence_theorem}
\end{equation}
\end{theorem}

\begin{proof}
Let $z(y)$ define the continuous inverse coordinate function, satisfying $y(z(y)) = y$. Since $\chi(z) \in C^2$, the coordinate density $\rho_{\text{clamp}}(z)$ is $C^1$-continuous on the compact domain $[0, H_{\text{SBL}}]$. Applying a second-order Taylor expansion to the inverse mapping $z(y)$ about the computational node $y_j$ yields:
\begin{equation}
    z(y_{j}) = z(y_{j-1}) + \frac{dz}{dy}(y_{j-1}) \Delta y + \frac{1}{2}\frac{d^2z}{dy^2}(\eta_j) (\Delta y)^2, \label{eq:taylor_expansion_inverse}
\end{equation}
where $\Delta y = \frac{1}{N_z - 1}$ and $\eta_j \in (y_{j-1}, y_j)$. Using the inverse derivative relations, we calculate:
\begin{equation}
    \frac{dz}{dy} = \frac{C_{\text{max}}}{\rho_{\text{clamp}}(z)}, \qquad \frac{d^2 z}{dy^2} = -\frac{C_{\text{max}}^2 \rho_{\text{clamp}}'(z)}{\left[ \rho_{\text{clamp}}(z) \right]^3}. \label{eq:inverse_derivs_calc}
\end{equation}
Because the clamped density is bounded away from zero ($\rho_{\text{clamp}} \ge \epsilon_z > 0$) and its derivatives are bounded, the second derivative is bounded pointwise by a positive constant $M > 0$. The global discretization error is bounded by:
\begin{equation}
    \max_{1 \le j \le N_z} \left| z_j - z(y_j) \right| \le \frac{1}{2} \max_{y \in [0,1]} \left| \frac{d^2 z}{dy^2} \right| (\Delta y)^2 \le \frac{M}{2 (N_z - 1)^2} = \mathcal{O}\left( N_z^{-2} \right), \label{eq:convergence_bound_step}
\end{equation}
completing the proof.
\end{proof}

\subsection{Fenichel Invariant Manifold Slaving and Stiffness Regularization}

Low-order Galerkin modal projections of non-linear boundary-layer budgets are prone to spurious, superlinear finite-time blowup ($t_{\text{blowup}} \le \frac{1}{2\sigma u_1(0)^2}$) near turning points, where the fast TKE eigenvalue vanishes ($\lambda_f \to 0$) \cite{mcnider1995}. This instability occurs because modal truncation allows the secondary fast mode $u_2$ to undergo negative excursions ($u_2 \propto -u_1^2$), introducing a positive cubic feedback term ($\sigma u_1^3$) into the primary slow mode equation \cite{mcnider1995}.

To regularize this stiffness, we apply \textit{Fenichel Invariant Manifold Slaving} \cite{fenichel1979}. We replace the prognostic ODE for the secondary mode, $\dot{u}_2$, with the regularized algebraic slow-manifold parameterization:
\begin{equation}
    u_2 = \Phi(u_1, S) \equiv -\frac{c_{\text{cross}} u_1^2}{2 \left| \lambda_f(S) \right| + \delta_m}, \label{eq:slaving_def}
\end{equation}
where $c_{\text{cross}} > 0$ is the spectral cross-coupling coefficient and $\delta_m \sim 10^{-3}\text{ s}^{-1}$ is a hyperbolic regularization parameter bounding the manifold solver near the fold locus where the fast eigenvalue vanishes \cite{fenichel1979}.

\begin{theorem}[Global Boundedness of Slaved Galerkin Trajectories]
\label{thm:fenichel_boundedness}
The algebraic substitution of the regularized Fenichel slaving relation \eqref{eq:slaving_def} into the primary slow mode equation eliminates the spurious cubic production feedback term $\sigma u_1^3$ entirely, establishing a globally bounded, stable attractor on the slow manifold $\mathcal{S}_\epsilon$.
\end{theorem}

\begin{proof}
Let the primary slow shear mode $u_1$ satisfy the governing modal ODE:
\begin{equation}
    \dot{u}_1 = G - c_{\text{shear}} u_1 - \beta u_1 u_2 - K_d u_1^3, \label{eq:u1_ode}
\end{equation}
where $G > 0$ is geostrophic forcing, $c_{\text{shear}} > 0$ is linear damping, $\beta > 0$ is coupling, and $K_d > 0$ is physical spectral dissipation. Substituting the Fenichel slaving relation \eqref{eq:slaving_def} into \eqref{eq:u1_ode} yields the closed system:
\begin{equation}
    \dot{u}_1 = G - c_{\text{shear}} u_1 + \left( \frac{\beta c_{\text{cross}}}{2 \left| \lambda_f(S) \right| + \delta_m} - K_d \right) u_1^3. \label{eq:u1_closed}
\end{equation}
Let $\mu(S) \equiv \frac{\beta c_{\text{cross}}}{2 \left| \lambda_f(S) \right| + \delta_m}$ represent the regularized production feedback rate. Because the positive regularization threshold $\delta_m > 0$ bounds the denominator from below, we have $\mu(S) \le \frac{\beta c_{\text{cross}}}{\delta_m} < \infty$ globally. By choosing a physical spectral dissipation parameter $K_d$ that satisfies the sub-critical dissipation constraint:
\begin{equation}
    K_d > \frac{\beta c_{\text{cross}}}{\delta_m}, \label{eq:dissipation_constraint}
\end{equation}
the effective cubic coefficient remains strictly negative for all shear states: $\Lambda(S) \equiv \mu(S) - K_d \le -\Lambda_0 < 0$. The dynamical system is bounded by the differential inequality $\dot{u}_1 \le G - c_{\text{shear}} u_1 - \Lambda_0 u_1^3$, defining a compact, invariant trapping region:
\begin{equation}
    \mathcal{K} = \left\{ u_1 \in \mathbb{R} : 0 \le u_1 \le \sqrt{\frac{G}{\Lambda_0}} + 1 \right\}. \label{eq:trapping_region}
\end{equation}
Any trajectory starting outside $\mathcal{K}$ has a strictly negative derivative ($\dot{u}_1 < 0$) and enters $\mathcal{K}$ in finite time, where it remains trapped for all future times ($t > 0$). This completes the proof.
\end{proof}

\subsection{Coupled Solver CFL Guarantees}

In production NWP applications, the maximum allowable integration step $\Delta t$ is governed by the Courant-Friedrichs-Lewy (CFL) limit of the vertical diffusion solver. Under our bounded grid-stretching and Fenichel slaving solver, we guarantee that the coupled solver remains stable under a non-vanishing minimum time step:

\begin{proposition}[Global CFL Stability Guarantee]
\label{prop:cfl_guarantee}
The maximum time step $\Delta t$ required to maintain numerical stability in the coupled TKE-shear solver remains strictly bounded from below by a non-zero physical limit, independent of the singular perturbation limit $\epsilon \to 0$ or the zero-shear limit $S \to 0$:
\begin{equation}
    \Delta t \ge \frac{(\Delta z_{\text{min}})^2}{2 K_{\text{max}}} \equiv \Delta t_{\text{CFL}} > 0, \label{eq:cfl_guarantee_formula}
\end{equation}
where $K_{\text{max}}$ is the maximum turbulent exchange coefficient bounded by the slow-manifold TKE trapping region.
\end{proposition}

\begin{proof}
An explicit vertical diffusion scheme requires $\Delta t \le \frac{(\Delta z_j)^2}{2 K_j}$ pointwise. By Theorem~\ref{thm:grid_bounds}, the discrete physical grid spacing is bounded from below by the clamped grid generator: $\Delta z_j \ge \Delta z_{\text{min}} > 0$. Furthermore, because Theorem~\ref{thm:fenichel_boundedness} establishes a globally bounded TKE ($E \le E_{\text{max}}$), the vertical turbulent exchange coefficients (which scale as $K = l \sqrt{E+\delta}$) are globally bounded from above by $K \le K_{\text{max}} < \infty$. This keeps the local explicit integration limit bounded away from zero: $\Delta t \ge \frac{(\Delta z_{\text{min}})^2}{2 K_{\text{max}}} > 0$, completing the proof.
\end{proof}